\subsection{1. Statement and answer}\label{statement-and-answer}

Write \(X=\{0,1\}^3\), ordered as \(000,001,010,011,100,101,110,111\).
Let

\[
\mathcal M=\left\{\lambda\bigotimes_{i=1}^3\mathrm{Ber}(u_i)
 +(1-\lambda)\bigotimes_{i=1}^3\mathrm{Ber}(v_i):
 \lambda,u_i,v_i\in[0,1]\right\}.
\]

This is the closure of binary RBM(3,1). Define

\[
\rho(p)=\min_{q\in\mathcal M}D(p\Vert q),\qquad
D(p\Vert q)=\sum_{x:p_x>0}p_x\log\frac{p_x}{q_x}.
\]

A positive target entry divided by zero gives infinite divergence. The
minimum exists: the parameter cube is compact, its image is compact,
divergence is lower semicontinuous, and the uniform distribution
supplies a finite competitor. Using the infimum over finite RBM
parameters gives the same value, by approximation of the mixture
parameters from the interior.

\textbf{Theorem (computer-assisted).}

\[
\boxed{\max_{p\in\Delta_7}\rho(p)
=c=-\frac34\log(2\sqrt3-3)
=\frac34\log\left(1+\frac2{\sqrt3}\right).}
\]

Numerically, \(c\approx0.575738814443\) nats, or
\(c/\log2\approx0.830615532444\) bits (12 decimal places). The only
maximizing target distributions are

\[
u_E=\tfrac14(\delta_{000}+\delta_{011}+\delta_{101}+\delta_{110}),
\qquad
u_O=\tfrac14(\delta_{001}+\delta_{010}+\delta_{100}+\delta_{111}).
\]

The proof has two parts. Section 2 gives an analytic lower bound and an
exact parity projection. Sections 3--6 give the universal upper bound
and classify all equality cases. The computer-assisted part is a finite
convex-cover lemma, not a sample or an asymptotic numerical estimate.

\subsection{2. Exact lower bound: the parity
projection}\label{exact-lower-bound-the-parity-projection}

Put \(s=\sqrt3\),

\[
a=\frac{2-s}{4},\qquad b=\frac{2s-3}{4},\qquad
q^*=(\tfrac14,a,a,b,a,b,b,3a).
\]

All entries are positive, and

\[
q^*=\alpha\delta_{000}+(1-\alpha)\mathrm{Ber}(t)^{\otimes3},
\qquad \alpha=\frac{3-s}{6},\quad t=\frac{3-s}{2}.
\tag{1}
\]

Thus \(q^*\in\mathcal M\). Direct substitution gives

\[
D(u_E\Vert q^*)=-\frac34\log(4b)=c.
\tag{2}
\]

It remains to exclude every better mixture; evaluating (2) alone would
not do that.

\subsubsection{2.1 Necessary supermodular
inequalities}\label{necessary-supermodular-inequalities}

For coordinatewise meet and join, let \(A\) have one row

\[
A_{xy}=e_{x\wedge y}+e_{x\vee y}-e_x-e_y
\]

for each incomparable unordered pair \(x,y\) in the three-cube. There
are nine rows. A positive probability table is log-supermodular in this
orientation when \(A\log q\ge0\).

Every strictly positive-parameter mixture in \(\mathcal M\) is
log-supermodular after independent coordinate complements. To see this
directly, complement coordinate \(i\) if necessary so that
\(u_i\le v_i\). Factoring out the first product gives

\[
\log q(x)=\text{a modular function of }x
 +\log(1+\exp(d+\sum_i w_i x_i)),\qquad w_i\ge0.
\]

The function \(\phi(z)=\log(1+e^z)\) has \(\phi''\ge0\). Its
supermodular difference is

\[
\phi(z+r+t)+\phi(z)-\phi(z+r)-\phi(z+t)
=\int_0^r\int_0^t\phi''(z+v+w)\,dw\,dv\ge0
\]

for \(r,t\ge0\). Hence all nine inequalities hold. Approximating
boundary mixture parameters proves the corresponding closed polynomial
inequalities there as well.

Complementing all three coordinates preserves the same cone, since it
exchanges meet with join. The eight orientations therefore have four
representatives. The even complements \(0,3,5,6\) represent all four and
preserve \(u_E\). Only this necessary inclusion is used; sufficiency of
the supermodularity conditions is not required by the proof.

\subsubsection{2.2 A three-multiplier convex
certificate}\label{a-three-multiplier-convex-certificate}

For arbitrary \(z\in\mathbb R^8\), define

\[
F(z)=\log\sum_xe^{z_x}-\sum_x(u_E)_x z_x.
\]

This is convex, and \(\nabla F(\log q^*)=q^*-u_E\). The three
square-face rows with their upper vertex at \(111\) are

\[
L_1=e_{001}+e_{111}-e_{011}-e_{101},
\] \[
L_2=e_{010}+e_{111}-e_{011}-e_{110},\qquad
L_3=e_{100}+e_{111}-e_{101}-e_{110}.
\]

They satisfy

\[
L_i\log q^*=0,\qquad
q^*-u_E=a(L_1+L_2+L_3),\qquad a>0.
\tag{3}
\]

Indeed, the active-face identities are \(b^2=3a^2\). The remaining six
supermodular inequalities are strict; they reduce to the positive
quantities \(b/4-a^2\) and \(3a/4-ab\).

For any \(z\) in this cone, the convex supporting-plane inequality and
(3) give

\[
F(z)\ge F(\log q^*)+a\sum_{i=1}^3L_i(z-\log q^*)
\ge F(\log q^*).
\]

Adding the constant \(\sum_x(u_E)_x\log(u_E)_x\) proves
\(D(u_E\Vert q)\ge c\) throughout this cone. Applying an even complement
gives the same bound in each of the other three orientations. Every
mixture with interior parameters belongs to one of these cones.
Approximation by interior mixture parameters extends the inequality to
all of \(\mathcal M\), including infinite-divergence cases. Together
with (1)--(2), this proves

\[
\rho(u_E)=\rho(u_O)=c.
\tag{4}
\]

For completeness, the four closest distributions to \(u_E\) are
precisely the even translates \(q^*(x\mathbin\oplus f)\),
\(f\in\{0,3,5,6\}\). For positive \(q\), the preceding supporting-plane
remainder equals \(D(q^*\Vert q)\); thus equality within a cone forces
\(q=q^*\). A boundary \(q\) with a zero entry cannot give equality:
along an interior approximation the same remainder tends to infinity
because \(q^*>0\). Odd translates give the four projections for \(u_O\).

The function \texttt{lower\_certificate()} in \texttt{verify.py} checks
(1), normalization, positivity, all nine polynomial inequalities, the
three active faces, (3), and the objective identity exactly in
\(\mathbb Q(\sqrt3)\).

\subsection{3. Reduce every target to six support
types}\label{reduce-every-target-to-six-support-types}

\subsubsection{3.1 Local-channel
monotonicity}\label{local-channel-monotonicity}

A stochastic map on one binary coordinate takes a Bernoulli product to
another Bernoulli product. It therefore preserves \(\mathcal M\), as
does a product of such maps. The log-sum inequality gives data
processing,

\[
D(Kp\Vert Kq)\le D(p\Vert q).
\]

Taking infima and using \(K\mathcal M\subseteq\mathcal M\) gives

\[
\rho(Kp)\le\rho(p).
\tag{5}
\]

Coordinate permutations and complements preserve \(\rho\) exactly, since
their inverses also preserve the model.

\subsubsection{3.2 Removing an entry by reversing a
channel}\label{removing-an-entry-by-reversing-a-channel}

Fix one coordinate, with two unnormalized slices \(u_y=p_{a,y}\) and
\(v_y=p_{b,y}\). Suppose both slices have nonempty support and

\[
\operatorname{supp}(u)\subseteq\operatorname{supp}(v).
\]

Define

\[
\tau=\min_{y:u_y>0}\frac{v_y}{u_y}>0,
\qquad r_{a,y}=(1+\tau)u_y,\quad r_{b,y}=v_y-\tau u_y.
\tag{6}
\]

The table \(r\) is nonnegative and has total mass one. It has no new
positive entry and has at least one fewer positive entry. Moreover,
\(p=Kr\), where

\[
K(a\mid a)=\frac1{1+\tau},\quad
K(b\mid a)=\frac\tau{1+\tau},\quad K(b\mid b)=1.
\tag{7}
\]

Consequently \(\rho(p)\le\rho(r)\) by (5). Comparable slices in the
other direction are treated by swapping \(a,b\).

If a coordinate has only one nonempty slice, the distribution belongs to
\(\mathcal M\): any two-bit distribution is a mixture of two products by
conditioning on either of its bits. Otherwise repeat (6) whenever
possible. In at most seven steps, the process reaches either that
zero-distance case or a support whose two projected slices are
incomparable in each coordinate.

\subsubsection{3.3 Exhaustive support
classification}\label{exhaustive-support-classification}

For a nonempty support \(S\subseteq\{0,\ldots,7\}\) and a bit
\(b\in\{1,2,4\}\), form

\[
S_0(b)=\{x\in S:x\mathbin\&b=0\},\qquad
S_1(b)=\{x\mathbin\oplus b:x\in S, x\mathbin\&b\ne0\}.
\]

Retain \(S\) if neither of these two sets contains the other, for all
three bits. Exhaustively applying this test to all 255 nonempty supports
retains 50. Their orbits under the 48 coordinate permutations and
complements have exactly the following representatives:

{
\begin{longtable}[]{@{}
  >{\raggedright\arraybackslash}p{(\linewidth - 6\tabcolsep) * \real{0.2000}}
  >{\raggedleft\arraybackslash}p{(\linewidth - 6\tabcolsep) * \real{0.2667}}
  >{\raggedleft\arraybackslash}p{(\linewidth - 6\tabcolsep) * \real{0.2667}}
  >{\raggedleft\arraybackslash}p{(\linewidth - 6\tabcolsep) * \real{0.2667}}@{}}
\toprule\noalign{}
\begin{minipage}[b]{\linewidth}\raggedright
Representative, in binary-index notation
\end{minipage} & \begin{minipage}[b]{\linewidth}\raggedleft
Orbit size
\end{minipage} & \begin{minipage}[b]{\linewidth}\raggedleft
Cover pieces
\end{minipage} & \begin{minipage}[b]{\linewidth}\raggedleft
Vertex checks
\end{minipage} \\
\midrule\noalign{}
\endhead
\bottomrule\noalign{}
\endlastfoot
\(0,7\) & 4 & 2 & 4 \\
\(0,3,5\) & 8 & 18 & 54 \\
\(0,1,2,7\) & 24 & 112 & 448 \\
\(0,3,5,6\) & 2 & 648 & 2,592 \\
\(0,1,2,4,7\) & 8 & 4,226 & 21,130 \\
\(0,1,2,5,6,7\) & 4 & 11,594 & 69,564 \\
\textbf{Total} & \textbf{50} & \textbf{16,600} & \textbf{93,792} \\
\end{longtable}
}

The support enumeration is separately implemented in
\texttt{support\_certificate()}, without importing the discovery code.
All weights on these supports remain arbitrary real probabilities.
Enumerating supports does not discretize those weights; the next section
covers their entire simplices. There are 52 prescribed roots in total
and maximum split depth 20; both statistics are reconstructed in the
replay receipt.

\subsection{4. The finite convex-cover
lemma}\label{the-finite-convex-cover-lemma}

\textbf{Lemma.} For each of the six support simplices above, there is a
finite simplicial cover such that every piece has a single explicit
witness \(q\in\mathcal M\) satisfying

\[
D(v\Vert q)\le c
\tag{8}
\]

at every vertex \(v\) of that piece. Every verified vertex inequality is
strict except when the vertex is exactly \(u_E\) or \(u_O\); those
equalities use the corresponding translated witness from Section 2.

This is the computer-assisted lemma. Its exact data are the six JSON
files in \texttt{certificates/}. Their hashes appear in
\texttt{verification.json} and the bundle manifest. The checker proves
the following finite assertions from those data; no discovery score is
accepted as evidence.

\subsubsection{4.1 The initial cover is
explicit}\label{the-initial-cover-is-explicit}

For four of the support types, the root is simply the full probability
simplex on the indicated vertices. For \(S=\{0,3,5,6\}\), take each of
the 24 orders \(x_1,x_2,x_3,x_4\) and the simplex with vertices

\[
r_j=\frac1j\sum_{i=1}^j\delta_{x_i},\qquad j=1,2,3,4.
\tag{9}
\]

These simplices cover the parity simplex. Indeed, order the weights so
that \(p_{x_1}\ge\cdots\ge p_{x_4}\), put \(p_{x_5}=0\), and use the
exact convex decomposition

\[
p=\sum_{j=1}^4 j(p_{x_j}-p_{x_{j+1}})r_j.
\tag{10}
\]

The coefficients are nonnegative and sum to one.

For \(S=\{0,1,2,4,7\}\), use all orders of the odd-parity vertices and
cone each simplex (9) from \(\delta_0\). Formula (10), with an
additional coefficient \(p_0\) at \(\delta_0\), proves coverage. All
these root vertices are rational and are reconstructed by the checker.

\subsubsection{4.2 Every subdivision preserves
coverage}\label{every-subdivision-preserves-coverage}

A split record names two distinct vertex positions \(i,j\). Let
\(m=(v_i+v_j)/2\). Replace the parent simplex by two children, replacing
\(v_i\) with \(m\) in one child and \(v_j\) with \(m\) in the other.

The children cover the parent: in barycentric coordinates \(\lambda\),
choose the first child if \(\lambda_i\le\lambda_j\), assigning weight
\(2\lambda_i\) to \(m\) and \(\lambda_j-\lambda_i\) to \(v_j\). The
other case is symmetric. Conversely both children lie in the parent. All
new vertices are rational.

Starting from the prescribed roots, the checker follows both children of
every split, reconstructs their vertices with exact fractions, and
requires every reached node to exist exactly once. It rejects invalid
split edges, missing children, duplicate addresses, and unreachable
extra nodes. Thus the final leaves form a genuine cover of every root
simplex.

\subsubsection{4.3 Witnesses belong to the requested
model}\label{witnesses-belong-to-the-requested-model}

There are only two witness formats.

\begin{enumerate}
\def\labelenumi{\arabic{enumi}.}
\tightlist
\item
  A \texttt{star} witness specifies an atom \(f\) and means the
  translated distribution \(q^*(x\mathbin\oplus f)\) from (1).
  Membership is exact in \(\mathbb Q(\sqrt3)\).
\item
  A \texttt{mixture} witness gives seven integers and a positive common
  denominator. They specify \(\lambda,u_1,u_2,u_3,v_1,v_2,v_3\). The
  checker verifies that all parameters are in \([0,1]\), reconstructs
  the eight probabilities as exact rational numbers, and checks
  normalization and nonnegativity.
\end{enumerate}

The search rounded mixture parameters, not probability-table entries.
Hence membership does not rest on a tolerance test for tensor rank, a
numerical decomposition, or approximate vanishing of a determinant.

\subsubsection{4.4 The logarithm inequalities use exact
arithmetic}\label{the-logarithm-inequalities-use-exact-arithmetic}

For a rational \(x>0\), first write \(x=2^k u\), \(1\le u<2\). Set
\(y=(u-1)/(u+1)\), so \(0\le y\le1/3\). Then

\[
\log u=2\sum_{j=0}^{N-1}\frac{y^{2j+1}}{2j+1}+R_N,
\quad
0\le R_N\le\frac{9}{4(2N+1)3^{2N+1}}.
\tag{11}
\]

The tail estimate follows by bounding all remaining denominators by
\(2N+1\) and summing a geometric series. The checker uses \(N=48\),
fixed-point scale \(2^{128}\), and integer floor/ceiling operations
after every multiplication and division. The same procedure bounds
\(\log2\). Negative \(k\) reverses the endpoints of the multiplied
interval. Thus every returned interval rigorously contains the
logarithm.

For algebraic witnesses, an integer square root gives

\[
\frac{L}{2^{128}}<\sqrt3<\frac{L+1}{2^{128}},
\qquad L=\left\lfloor\sqrt{3\,2^{256}}\right\rfloor.
\]

The checker verifies both bounding square inequalities. Affine interval
evaluation and monotonicity of the logarithm then enclose every star
probability and the constant \(c\).

For each non-equality vertex, the checker evaluates an upper bound on

\[
\sum_xv_x\log v_x-\sum_xv_x\log q_x
\]

and proves it is strictly below a lower bound on \(c\). A zero target
term is omitted; a positive target over a zero witness entry is
rejected. No floating-point logarithm or floating-point comparison
participates in these proof decisions.

At an equality vertex, the checker requires the exact uniform parity
vector and a star atom of matching parity. The equality is then the
exact identity (2), not a numerical tolerance exception. All other
vertices require a strict inequality. The smallest certified strict
vertex slack in this cover exceeds \(2\times10^{-6}\) nats. This is a
margin for the finite vertex tests, not a uniform separation of all
nonmaximizing targets from \(c\).

\subsubsection{4.5 Replay result}\label{replay-result}

Run, using Python 3.10 or later with only its standard library:

\begin{verbatim}
python3 verify.py --out replay.json
python3 test_verify.py
\end{verbatim}

The first command reconstructs the lower-bound algebra, all support
orbits, all six covers, and all 93,792 vertex inequalities. The recorded
local run returned \texttt{PASS}, with 16,600 leaves. There are 552
incidences of exact parity vertices across overlapping pieces; every
other vertex test is strict. The second command passed 13 tests,
including malformed-certificate rejection, an independent
Decimal-logarithm diagnostic, and 765 exact rational
channel-reconstruction checks.

The verifier imports no optimization package or discovery
implementation. Numerical search suggested the witnesses; the exact
replay proves the finite lemma. The remaining trusted base is the
mathematical argument in this document, the checker implementation,
Python's standard-library integer/fraction semantics, and the execution
platform. No claim of proof-assistant certification or independent
specialist review is made.

\subsection{5. Universal upper bound}\label{universal-upper-bound}

Fix a support-simplex point \(p\). By Section 4 it belongs to a leaf
simplex with vertices \(v_1,\ldots,v_m\) and a common witness
\(q\in\mathcal M\). Write \(p=\sum_j\lambda_jv_j\). For fixed \(q\),
divergence is convex in its first argument, since \(t\log t\) is convex
and the cross-entropy term is linear. Therefore

\[
\rho(p)\le D(p\Vert q)
\le\sum_j\lambda_jD(v_j\Vert q)\le c.
\tag{12}
\]

This establishes the bound for every real weight vector on all six
representative supports. Cube symmetries give the same conclusion on
every irreducible support. Reversing the reductions in Section 3 and
applying (5) gives \(\rho(p)\le c\) for every distribution on all eight
cube vertices. Together with (4), this proves the exact global value.

Notice the quantifier: (12) covers every point in each simplex, not just
its vertices as a finite list of target distributions. Convexity is
precisely what turns the finite certificate into a universal upper
bound.

\subsection{6. Classification of all maximizing
targets}\label{classification-of-all-maximizing-targets}

First consider a target on an irreducible support. Every non-parity
vertex of every leaf has a strict bound. If a target is not a leaf
vertex, its barycentric representation uses at least two distinct
vertices. On the finite-divergence simplex, \(D(\cdot\Vert q)\) is
strictly convex: \(\sum_xp_x\log p_x\) is strictly convex and the other
term is linear. Thus (12) is strict for any such target. It follows that
equality on irreducible supports is possible only at a uniform parity
vector. Cube symmetries only exchange the two parity vectors.

Now take a target for which Section 3 made at least one strict support
reduction. If its terminal target is not uniform parity, monotonicity
already gives \(\rho(p)<c\). Otherwise consider the last inverse channel
\(K\) in (7), so its input is uniform parity \(u\). Let \(q^*_u>0\) be a
parity projection from Section 2. Then

\[
\rho(Ku)\le D(Ku\Vert Kq^*_u)<D(u\Vert q^*_u)=c.
\tag{13}
\]

The middle inequality is strict by the equality condition in the log-sum
inequality. In each output fiber with changed bit \(b\), the channel
merges the two input values \(a,b\) with positive \(q^*_u\)-mass from
both. Exactly one of the two parity target entries is positive. Their
likelihood ratios relative to \(q^*_u\) are consequently different,
while \(\tau>0\) ensures genuine merging. Earlier channels preserve the
strict upper bound by (5).

Targets with a deterministic coordinate have zero error. Hence the only
global maximizers are \(u_E\) and \(u_O\), as claimed.

\subsection{7. Constructive content and dimensional
transfer}\label{constructive-content-and-dimensional-transfer}

The upper bound supplies an explicit procedure, not just an existence
argument. For a given target, apply (6), recording each channel.
Transform the terminal support to a representative. Sort its weights to
choose a root in (9), if needed. Descend the split tree using
barycentric coordinates. Read the exact witness at its leaf, undo the
cube symmetry, and apply the recorded channels. The resulting
distribution remains a mixture of two products and has divergence at
most \(c\).

The command-line implementation \texttt{witness.py} accepts exact
rational input probabilities. It verifies all six certificates before
answering queries, performs the channel and barycentric calculations
with fractions, and returns mixture weights and Bernoulli parameters in
\(\mathbb Q(\sqrt3)\), together with the reduction, symmetry and leaf
trace. It returns a feasible approximation certified to have divergence
at most \(c\), \textbf{not} a generally optimal projection or an
algorithm for training an RBM. The written procedure applies to
arbitrary real targets; the executable interface is restricted to
rational targets. Bit-coordinate arrays in this interface are explicitly
least-significant-bit first.

\subsubsection{7.1 Worked channel and tree
query}\label{worked-channel-and-tree-query}

Take \(p=(3,1,0,4,0,4,3,1)/16\). For the least significant bit, the
slices with bit values 0 and 1 have comparable supports, and (6) gives
\(\tau=1/3\). The reduced target is exactly \(u_E\). The deterministic
tie rule orders its support as \((0,3,5,6)\), selecting root 0 with
barycentric coordinates \((0,0,0,1)\). Descending the recorded tree
locates a leaf containing that parity vertex and its matching star
witness. Undoing the recorded channel produces a two-product mixture for
\(p\). The strict data-processing argument in (13) shows its error is
below \(c\). The exact leaf address, final parameters and barycentric
weights are retained in \texttt{examples/channel-witness.json}.

\begin{verbatim}
python3 witness.py 3/16 1/16 0 1/4 0 1/4 3/16 1/16 --out query.json
python3 test_witness.py
\end{verbatim}

\subsubsection{7.2 A conditioning transfer
theorem}\label{a-conditioning-transfer-theorem}

For positive integers \(n,k\), let \(\mathcal M_{n,k}\) denote mixtures
of at most \(k\) Bernoulli products on \(n\) bits, with closed
parameters, and let
\(R(n,k)=\sup_p\min_{q\in\mathcal M_{n,k}}D(p\Vert q)\).

\textbf{Proposition (conditioning lift).} If \(n\ge m\ge1\), then

\[
R(n,2^{n-m}k)\le R(m,k).
\tag{14}
\]

\textbf{Proof.} Write a target as \(p(z,y)=p_Z(z)p(y\mid z)\), with
\(z\) consisting of \(n-m\) fixed coordinates. Ignore zero-mass
contexts. For each other context choose a \(k\)-product approximation
\(q_z\) to its conditional law with error at most \(R(m,k)\). The
distribution \(q(z,y)=p_Z(z)q_z(y)\) is a mixture of at most
\(2^{n-m}k\) products: append the deterministic coordinates \(z\) to
every component of \(q_z\), and multiply its weight by \(p_Z(z)\). The
supports of the contexts are disjoint, so cancellation of the common
context weights gives

\[
D(p\Vert q)=\sum_{z:p_Z(z)>0}p_Z(z)D(p(\cdot\mid z)\Vert q_z)
\le R(m,k).
\]

Taking the minimum and supremum proves (14). This is the standard
disjoint-support mixture principle, written here explicitly; no novelty
is claimed for the lifting mechanism.

\textbf{Corollary.} The three-bit theorem implies, for every \(n\ge3\)
and every \(k\ge2^{n-2}\),

\[
R(n,k)\le c=-\frac34\log(2\sqrt3-3).
\tag{15}
\]

In bits this upper bound is \(c/\log2=0.830615532444310205\ldots\). It
improves the elementary one-bit upper bound at \(k=2^{n-2}\):
conditioning on \(n-2\) bits and approximating each remaining two-bit
table by its independent marginals gives error \(I(X;Y)\le\log2\). The
mechanism is established in the disjoint-support mixture literature; the
quantitative improvement is a consequence of the candidate three-bit
constant. Sharpness for \(n>3\), maximizers in higher dimensions,
stronger bounds at other component budgets, and consequences for general
multi-hidden-unit RBMs are \textbf{not} established. In particular,
product mixtures and RBMs must not be interchanged for these larger
parameter values.

\texttt{Oracle.approximate} implements (14) for \(m=3,k=2\). Its tests
cover 765 rational three-bit targets over every nonempty support, both
parity targets, the worked channel example, and 50 four-/five-bit
conditioning examples. These finite tests diagnose the implementation;
the all-dimensional claim is the analytic argument above.

The retained discovery scripts are optional. They are not imported by
the checker and need NumPy/SciPy only if the search is rerun. A
numerical run need not regenerate byte-identical witnesses on a
different optimizer version; the frozen exact certificate, not
deterministic discovery, is the proof object.

No boundary-MLE stratum classification, floating-point EM fit,
approximate global optimizer, unproved entropy inequality, or symmetry
average of mixtures is needed in the final proof.

\subsection{8. Prior art, scope, and review
boundary}\label{prior-art-scope-and-review-boundary}

The problem is the frozen AIM/UnsolvedMath item
\texttt{AIM-PROBABILITY-0002}, not the neighboring
RBM(3,2)/three-component-mixture problem. Montúfar's 2018 review,
Section 9, item 10, already records the proposed value, attributed to
discussions with Johannes Rauh. Its formulation and numerical constant
are not claimed as discoveries here. The current full-scope contribution
proposed for review is the channel reduction plus the exact convex
cover, together with the short parity certificate and equality argument.

The necessary supermodularity description is consistent with
Allman--Rhodes--Sturmfels--Zwiernik's nonnegative-rank-two framework.
Section 2 proves the inclusion actually needed, so the final argument
does not silently rely on a positive-support characterization at zero
entries. The exact boundary-MLE paper was useful during reconnaissance
but is not a dependency of the final proof.

Alexandr and Hoşten (2025), Proposition 2, use strict convexity to
localize divergence maximizers at vertices of logarithmic Voronoi
polytopes for linear and toric models. The shared convexity principle is
established material. Our cover consists instead of prescribed simplices
with feasible common witnesses for a nonconvex two-product mixture
model: a witness need not be the projection throughout its simplex. The
new proof object does not reconstruct logarithmic Voronoi polytopes.
Montúfar, Rauh and Ay (2013), Lemma 6, give the established
disjoint-support decomposition underlying Section 7.2.

The bounded novelty searches found no source proving this same global
equality. That does not establish priority, present-day openness, or
absence of unpublished work. The accompanying novelty note records the
search limits. Internal model-mediated review and producer replay are
not independent mathematical review. The constructive interface and the
conditioning corollary demonstrate scoped reuse, not measured impact or
a substantiated four-star assessment.

\subsection{9. Proof-to-checker map and
disclosure}\label{proof-to-checker-map-and-disclosure}

{
\begin{longtable}[]{@{}
  >{\raggedright\arraybackslash}p{(\linewidth - 4\tabcolsep) * \real{0.3333}}
  >{\raggedright\arraybackslash}p{(\linewidth - 4\tabcolsep) * \real{0.3333}}
  >{\raggedright\arraybackslash}p{(\linewidth - 4\tabcolsep) * \real{0.3333}}@{}}
\toprule\noalign{}
\begin{minipage}[b]{\linewidth}\raggedright
Obligation
\end{minipage} & \begin{minipage}[b]{\linewidth}\raggedright
Exact executable check
\end{minipage} & \begin{minipage}[b]{\linewidth}\raggedright
Analytic bridge retained in the manuscript
\end{minipage} \\
\midrule\noalign{}
\endhead
\bottomrule\noalign{}
\endlastfoot
Parity lower bound & \texttt{lower\_certificate}: mixture identity,
signs, nine inequalities, three multipliers & Necessary cone inclusion,
convex supporting plane, boundary passage and uniqueness; Section 2 \\
All targets reduce to six support orbits &
\texttt{support\_certificate}: all 255 supports, 50 irreducibles and six
orbits & Reverse-channel formula, strict support decrease and data
processing; Section 3 \\
Continuous simplex coverage & \texttt{root\_simplices},
\texttt{verify\_cover}: roots, every branch, no orphan node &
Ordered-weight root decomposition and midpoint coverage; Sections
4.1--4.2 \\
Feasible common witnesses & \texttt{witness\_logs},
\texttt{lower\_certificate} & Product-mixture model definition; Sections
2 and 4.3 \\
Strict vertex error bounds & \texttt{log\_interval},
\texttt{vertex\_slack}, exact parity exception & Log-series remainder
and first-argument convexity; Sections 4.4--5 \\
Only two maximizers & Strict vertex tests and exact parity recognition &
Strict convexity and strict log-sum equality condition; Section 6 \\
Constructive query and lift & \texttt{witness.py},
\texttt{test\_witness.py} & Channel reversal and disjoint-context KL
identity; Section 7 \\
\end{longtable}
}

The checker is a standalone implementation separated from discovery
within the same producer-coordinated workflow. The test suite is not a
formal verification. Neither the lower-bound analytic bridges nor the
all-dimensional quantifiers are reduced to checking finitely many random
inputs.

The scholarly creator is Anonymous; Evidence Press is the publisher, not
an inferred personal author. A human directed the research and supplied
review; an OpenAI Codex research assistant drafted the mathematical
argument, implemented search and verification, revised the manuscript
and prepared the release. Producer-coordinated model reports provide
internal editorial scrutiny only. No external funding or competing
interests have been declared; no human/animal subjects or personal
datasets are involved. These disclosures do not establish independent
validation.

All original prose and mathematical certificate data are offered under
CC0-1.0; original code is MIT. The package links third-party literature
and code without relicensing them. The supplied review is not
redistributed. The paper, source, exact certificates, executable witness
interface, semantic controls and original response matrix are available
together in the versioned evidence package. Optional discovery requires
the separately listed numerical dependencies; all proof replay and query
operations use Python's standard library.

\subsubsection{Sources}\label{sources}

\begin{enumerate}
\def\labelenumi{\arabic{enumi}.}
\tightlist
\item
  Guido Montúfar, \emph{Restricted Boltzmann Machines: Introduction and
  Review} (2018), especially Sections 3 and 9, item 10.
  \href{https://arxiv.org/abs/1806.07066}{Primary preprint}.
\item
  Elizabeth S. Allman, John A. Rhodes, Bernd Sturmfels, Piotr Zwiernik,
  \emph{Tensors of Nonnegative Rank Two}, Linear Algebra and its
  Applications 473 (2015), 37--53.
  \href{https://arxiv.org/abs/1305.0539}{Primary preprint}.
\item
  Anna Seigal and Guido Montúfar, \emph{Mixtures and products in two
  graphical models}. This is the neighboring, different model.
  \href{https://arxiv.org/abs/1709.05276}{Primary preprint}.
\item
  Elizabeth S. Allman, Hector Baños Cervantes, Robin Evans, Serkan
  Hoşten, Kaie Kubjas, Daniel Lemke, John A. Rhodes and Piotr Zwiernik,
  \emph{Maximum Likelihood Estimation of the Latent Class Model through
  Model Boundary Decomposition}, Journal of Algebraic Statistics 10(1)
  (2019), 51--84. \href{https://doi.org/10.18409/jas.v10i1.75}{DOI};
  \href{https://arxiv.org/abs/1710.01696}{primary preprint}.
  Reconnaissance only.
\item
  Yulia Alexandr and Serkan Hoşten, \emph{Maximum information divergence
  from linear and toric models}, Information Geometry 8 (2025),
  159--197. \href{https://doi.org/10.1007/s41884-025-00166-3}{DOI};
  \href{https://arxiv.org/abs/2308.15598}{primary preprint}, Proposition
  2.
\item
  Guido Montúfar, Johannes Rauh and Nihat Ay, \emph{Maximal Information
  Divergence from Statistical Models Defined by Neural Networks}, in
  Geometric Science of Information (2013), 759--766.
  \href{https://arxiv.org/abs/1303.0268}{Primary preprint}, Lemma 6.
  Established disjoint-support mixture mechanism.
\end{enumerate}

The primary sources, GitHub implementation, and live Evidence Press
context are identified in \texttt{NOVELTY.md}. Original screening
excerpts remain in the preserved review-v1 archive and are not
redistributed in this successor.
